\section{Discussion}
\label{sec:discussion}

\subsection{Operational evidence is a digital-laboratory object}

The analysis shows that many laboratory-automation failures arise when a digital representation is treated as operationally complete without the evidence required to support that status. A listed driver does not establish accessible control. A method file does not identify a deployment. A geometry value does not establish physical portability. A telemetry stream does not reconstruct material state. A passing simulation does not establish wet or assay performance. These gaps are representational, although their consequences are physical.

The appropriate response is to make applicability, provenance, version, uncertainty, validation stage, and ownership explicit in the data model. This principle underlies the field-level metrics and the released schemas. It also explains why unknown remains distinct from absent, why source-reported evidence remains labeled, and why final material disposition is represented separately from software continuation. The resulting objects can be queried and validated without implying that the ordinal scores are calibrated probabilities.

\subsection{Observability should connect commands, materials, and outcomes}

Observability and recovery form the strongest technical relationship in the pilot. The association remains stable under every leave-one-thread-out deletion, and incident episodes provide a direct mechanism. Once physical actions have executed, diagnosis and recovery depend on knowing which commands completed, which material and resource identities were affected, which observations were available, which interventions occurred, and whether modeled and physical state were reconciled.

The run-event schema released with the Observatory represents this chain through stable identifiers and typed events. Its prospective evaluation should test whether greater linkage completeness predicts time to a verified safe state, unresolved-state count, operator burden, salvage, final disposition completeness, and incident recurrence. This moves the mechanism from forum-derived evidence to an operational dataset with incident and run denominators.

\subsection{Digital resources require evidence grades and maintenance state}

Physical-resource and resolved-knowledge records share a common stewardship problem. A field can be declared, measured, validated on one device, or independently reproduced. An answer can be active, review due, disputed, superseded, or archived. Collapsing these states creates apparent completeness at the cost of operational credibility.

The resource schemas therefore bind values to evidence grade, applicability, date, maintainer, known limitations, and lineage. This structure makes staleness and disagreement visible. It also permits correction without silently rewriting a previous release. The strongest positive cases in the corpus followed this pattern: local failure was converted into public diagnosis, an artifact was corrected, and the correction remained linked to its provenance.

\subsection{Robustness analyses reveal which results deserve priority}

Alternative partial-score weights show that the component metrics depend on how partial evidence is represented. Their central values are useful rubric summaries and should not be read as probabilities. Leave-one-thread-out analysis distinguishes relationships that remain stable across the purposive register from relationships that are concentrated in a small stratum. Observability--recovery, observability--validation, and interface accessibility--scheduling remain the strongest prospective priorities. The AI--validation relationship remains substantively important and sample-sensitive.

Alternative denominator analyses serve a related function. They expose which fields define the operational construct and how narrower subsets change the summary. The headline denominator remains the declared study object; alternative denominators are reported as adversarial checks, not selected after seeing a more favorable value. This pattern is suitable for digital-laboratory research where small, heterogeneous case sets are necessary during early metric construction.

\subsection{From public discussion to reusable infrastructure}

The forum performs distributed requirements elicitation and repair. Replies reveal missing constraints, distinguish evidence stages, connect users to hidden documentation, and occasionally convert an incident into maintained code or guidance. The proposed infrastructure preserves those contributions without replacing the source discussion. A canonical record points to the forum, maintained artifact, evidence grade, applicable versions, known limitations, and correction status.

The value of this infrastructure must be measured prospectively. Question templates can be evaluated through clarification turns and time to actionable response. Knowledge records can be evaluated through recurrence, staleness, correction latency, and links to maintained artifacts. Interface records can be evaluated through integration time and failed dispatches. Event schemas can be evaluated through diagnosis and recovery outcomes. These measures create denominators unavailable in the public discussion corpus.

\subsection{Implications for AI-enabled digital laboratories}

AI method generation inherits every unresolved representation in the laboratory stack. A model cannot safely infer a licence condition, missing geometry, unrecorded calibration, unstated stability window, or physical outcome that the system does not represent. Evaluation should therefore include context elicitation, uncertainty handling, abstention on underdetermined actions, human correction burden, unsafe omission, and performance at each physical validation stage.

The Context Expansion Ratio provides an early measure of this requirement. The selected discussion moved from a small action vocabulary to a materially larger execution and deployment context after community review. Future benchmarks should encode missing-context cases and score whether systems request the information needed to make a physical action and its claim scope defensible.
